% eksperymentalne tłumaczenie na włoski

%

\section{Potęga punktu względem okręgu} % lang-pl
\index{potęga punktu względem okręgu}% % lang-pl
\section{Potenza di un punto rispetto a una circonferenza} % lang-it
\index{potenza di un punto rispetto a una circonferenza}% % lang-it

Zdefiniujemy potęgę punktu względem okręgu, a potem potęgową oś i środek. % lang-pl
\todofoot{Inwersja: Audin \cite[s. 84]{audin_2003}.} % lang-pl
Definiremo la potenza di un punto rispetto a una circonferenza e successivamente l'asse radicale e il centro radicale. % lang-it
\todofoot{Inversione: Audin \cite[p. 84]{audin_2003}.} % lang-it
% zachowuje kąty
\todofoot{Coxeter s. 85, coaxial circles}

\begin{definition}
[potęga punktu względem okręgu] % lang-pl
[potenza di un punto rispetto a una circonferenza] % lang-it
	\label{def_power_point}
	Dane są okrąg $\Gamma$ o środku $O$ i promieniu $r$ oraz dowolny punkt $A$. % lang-pl
	Liczbę rzeczywistą % lang-pl
	Siano dati una circonferenza $\Gamma$ di centro $O$ e raggio $r$, nonché un punto arbitrario $A$. % lang-it
	Il numero reale % lang-it
	\begin{equation}
		\Pi(A) := |OA|^2 - r^2
	\end{equation}
	nazywamy potęgą punktu $A$ względem okręgu $\Gamma$. % lang-pl
	si chiama potenza del punto $A$ rispetto alla circonferenza $\Gamma$. % lang-it
\end{definition}

Pisze o tym Audin \cite[s. 89]{audin_2003}, nieznany autor w $\Delta_{84}^{11}$, Joanna Jaszuńska w $\Delta_{12}^2$ i $\Delta_{12}^3$, Bartłomiej Bzdęga w~$\Delta_{19}^{11}$. % lang-pl
Ne scrivono Audin \cite[p. 89]{audin_2003}, un autore anonimo in $\Delta_{84}^{11}$, Joanna Jaszuńska in $\Delta_{12}^2$ e $\Delta_{12}^3$, e Bartłomiej Bzdęga in $\Delta_{19}^{11}$. % lang-it
% https://www.deltami.edu.pl/2012/02/potega-punktu/
% https://www.deltami.edu.pl/2012/03/osie-potegowe/
% https://www.deltami.edu.pl/2019/11/potega-punktu-wzgledem-okregu/

O czymś na kształt potęgi punktu napisze Louis Gaultier w 1813 roku, ale sam termin (,,potęga punktu'') wprowadzi Jakob Steiner w 1826 roku. % lang-pl
Louis Gaultier scriverà di un concetto simile alla potenza di un punto nel 1813, ma il termine stesso (``potenza di un punto'') sarà introdotto da Jakob Steiner nel 1826. % lang-it
% Gaultier: published in the Journal de L'Lcole Polyttcknique, Cahier 16 (1813), pp. 124-214
\index[persons]{Steiner, Jakob}%
\index[persons]{Gaultier, Louis}%
Użyje jej, aby znaleźć okrąg przecinający cztery dane okręgi pod tym samym kątem; rozwiązać problem Apoloniusza (problem \ref{problem_apoloniusza}: znaleźć okrąg styczny do trzech); skonstruować okręgi Malfattiego (problem \ref{malfatti_problem}: trzy okręgi, które są styczne do pozostałych dwóch oraz boków zadanego trójkąta). % lang-pl
La utilizzerà per trovare una circonferenza che intersechi quattro circonferenze date sotto lo stesso angolo; risolvere il problema di Apollonio (problema \ref{problem_apoloniusza}: trovare una circonferenza tangente a tre date); costruire le circonferenze di Malfatti (problema \ref{malfatti_problem}: tre circonferenze tangenti alle altre due e ai lati di un triangolo assegnato). % lang-it
% TODO: index, ref
\index{problem!Apoloniusza}% % lang-pl
\index{okrąg!Malfattiego}% % lang-pl
Wydaje się, że ta lista jest niekompletna. % lang-pl
\index{problema!di Apollonio}% % lang-it
\index{circonferenza!di Malfatti}% % lang-it
Sembra che questo elenco non sia completo. % lang-it

Wśród elementarnych zastosowań można wskazać dowód twierdzenia o przecinających się cięciwach (fakt \ref{prop_intersecting_chords}). % lang-pl
Twierdzenie Pitagorasa dostarcza prostej interpretacji wielkości $\Pi(A)$, gdy punkt $A$ leży na zewnątrz okręgu: jest to długość stycznej do okręgu, która przechodzi przez punkt. % lang-pl
\index{twierdzenie!Pitagorasa} % lang-pl
Tra le applicazioni elementari si può citare la dimostrazione del teorema delle corde che si intersecano (fatto \ref{prop_intersecting_chords}). % lang-it
Il teorema di Pitagora fornisce una semplice interpretazione della quantità $\Pi(A)$ quando il punto $A$ si trova all'esterno della circonferenza: essa è la lunghezza della tangente alla circonferenza passante per il punto. % lang-it
\index{teorema!di Pitagora} % lang-it

\begin{proposition}
	Środki odcinków wspólnych stycznych do dwóch okręgów leżą na jednej prostej. % lang-pl
	I punti medi dei segmenti determinati dalle tangenti comuni a due circonferenze giacciono su una stessa retta. % lang-it
\end{proposition}

Nie wiemy, jak trudny jest dowód powyższego bez poniższego. % lang-pl
Non sappiamo quanto sia difficile dimostrare il risultato precedente senza utilizzare quello seguente. % lang-it

\begin{proposition}
[oś potęgowa istnieje] % lang-pl
[esistenza dell'asse radicale] % lang-it
\label{guzicki_6_11}%
    Dane są dwa niewspółśrodkowe okręgi $\Gamma_1(O_1, r_1)$ oraz $\Gamma_2(O_2, r_2)$. % lang-pl
    Wtedy miejscem geometrycznym punktów $P$ mających równe potęgi względem obu okręgów: % lang-pl
    Siano date due circonferenze non concentriche $\Gamma_1(O_1, r_1)$ e $\Gamma_2(O_2, r_2)$. % lang-it
    Allora il luogo geometrico dei punti $P$ aventi la stessa potenza rispetto a entrambe le circonferenze: % lang-it
	\begin{equation}
		\{S : \Pi(S, \Gamma_1) = \Pi(S, \Gamma_2)\}
	\end{equation}
	jest prosta prostopadła do prostej przechodzącej przez środki obu okręgów. % lang-pl
	Nazywamy ją \emph{osią potęgową} okręgów $\Gamma_1$, $\Gamma_2$. % lang-pl
	\index{oś potęgowa}% % lang-pl
	\index{potęgowa!oś|see{oś potęgowa}} % lang-pl
	è una retta perpendicolare alla retta passante per i centri delle due circonferenze. % lang-it
	La chiamiamo \emph{asse radicale} delle circonferenze $\Gamma_1$, $\Gamma_2$. % lang-it
	\index{asse radicale}% % lang-it
	\index{radicale!asse|see{asse radicale}} % lang-it
\end{proposition}

O osi potęgowej piszą Guzicki \cite[s. 173, 174]{guzicki_2021}, Bogdańska, Neugebauer \cite[s. 69]{neugebauer_2018}, Audin \cite[s. 89]{audin_2003}, Eves \cite[s. 92]{eves1_1972}; patrz też \cite[s. 11]{komisarczyk_2024}. % lang-pl
Oś potęgowa okręgów stycznych to prosta styczna do obydwu okręgów; oś potęgowa okręgów, które się przecinają w dwóch punktach, to prosta przez punkty przecięcia (to ćwiczenie u Hartshorne'a \cite[s. 182]{hartshorne_2000}). % lang-pl
Dell'asse radicale scrivono Guzicki \cite[p. 173, 174]{guzicki_2021}, Bogdańska e Neugebauer \cite[p. 69]{neugebauer_2018}, Audin \cite[p. 89]{audin_2003}, Eves \cite[p. 92]{eves1_1972}; si veda anche \cite[p. 11]{komisarczyk_2024}. % lang-it
L'asse radicale di due circonferenze tangenti è la loro tangente comune; l'asse radicale di due circonferenze che si intersecano in due punti è la retta passante per tali punti di intersezione (si tratta di un esercizio in Hartshorne \cite[p. 182]{hartshorne_2000}). % lang-it

\begin{proposition}
[środek potęgowy] % lang-pl
[centro radicale] % lang-it
	Dane są trzy parami niewspółśrodkowe okręgi $\Gamma_1, \Gamma_2, \Gamma_3$. % lang-pl
	Wtedy albo środki tych okręgów są współliniowe (i trzy osie potęgowe każdej pary są równoległe), albo wszystkie trzy osie przecinają się w jednym punkcie: \emph{środku potęgowym} okręgów. % lang-pl
	\index{środek potęgowy}% % lang-pl
	\index{potęgowy!środek|see{środek potęgowy}} % lang-pl
	Siano date tre circonferenze $\Gamma_1, \Gamma_2, \Gamma_3$ a due a due non concentriche. % lang-it
	Allora o i centri delle tre circonferenze sono allineati (e i tre assi radicali delle coppie di circonferenze sono paralleli), oppure tutti e tre gli assi si incontrano in un unico punto, detto \emph{centro radicale} delle circonferenze. % lang-it
	\index{centro radicale}% % lang-it
	\index{radicale!centro|see{centro radicale}} % lang-it
\end{proposition}

Terminy ,,oś potęgowa'' i ,,środek potęgowy'' to również pomysł wspomnianego Gaultiera. % lang-pl
Anche i termini ``asse radicale'' e ``centro radicale'' si devono al già citato Gaultier. % lang-it

Nie jest jasne, dlaczego niektórzy nazywają to twierdzeniem Monge'a (jak Bogdańska, Neugebauer \cite[s. 70]{neugebauer_2018}) albo też rozwiązaniem problemu Monge'a (jak Dörrie \cite[s. 151]{dorrie_1965}). % lang-pl
\index{twierdzenie!Monge'a}% % lang-pl
Ale jest jasne, że pozwala łatwo skonstruować oś potęgową dwóch okręgów: wystarczy przeciąć je jednocześnie trzecim okręgiem na dwa różne sposoby. % lang-pl
Non è chiaro perché alcuni chiamino questo risultato teorema di Monge (come Bogdańska e Neugebauer \cite[p. 70]{neugebauer_2018}) oppure soluzione del problema di Monge (come Dörrie \cite[p. 151]{dorrie_1965}). % lang-it
\index{teorema!di Monge}% % lang-it
È però chiaro che esso permette di costruire facilmente l'asse radicale di due circonferenze: basta intersecarle entrambe con una terza circonferenza in due modi differenti. % lang-it
% TODO: dla innych to jest twierdzenie, że Theorem 1.1 (Monge’s theorem of three circles) For any three circles in a plane, none of which is inside one of the others, the intersection points of each of the three pairs of external tangent lines are collinear. jak w https://atcm.mathandtech.org/EP2016/contributed/4052016_21160.pdf

Inną konstrukcję opisuje Hartshorne \cite[s. 182]{hartshorne_2000}: dane są dwa okręgi o środkach w punktach $O_1, O_2$ oraz promieniach odpowiednio $O_1A_1$, $O_2A_2$. % lang-pl
Niech $B$ będzie środkiem $A_1A_2$, zaś $P$ punktem przecięcia się prostych: prostopadłej do $O_1B$ przez $A_1$ oraz prostopadłej do $O_2B$ przez $A_2$. % lang-pl
Wtedy prostopadła do $O_1O_2$ przez $P$ jest szukaną osią potęgową. % lang-pl
Hartshorne \cite[p. 182]{hartshorne_2000} descrive un'altra costruzione: siano date due circonferenze con centri nei punti $O_1$, $O_2$ e raggi rispettivamente $O_1A_1$, $O_2A_2$. % lang-it
Sia $B$ il punto medio di $A_1A_2$ e sia $P$ il punto d'intersezione delle rette: la perpendicolare a $O_1B$ passante per $A_1$ e la perpendicolare a $O_2B$ passante per $A_2$. % lang-it
Allora la perpendicolare a $O_1O_2$ passante per $P$ è l'asse radicale cercato. % lang-it

Patrz też Guzicki \cite[s. 174]{guzicki_2021}. % lang-pl
Si veda anche Guzicki \cite[p. 174]{guzicki_2021}. % lang-it

% Neugebauer s. 71
\begin{proposition}
	Ortocentrum trójkąta jest środkiem potęgowym rodziny wszystkich okręgów o średnicach będących czewianami (odcinkami łączącymi wierzchołek z przeciwległym bokiem). % lang-pl
	L'ortocentro di un triangolo è il centro radicale della famiglia di tutte le circonferenze aventi come diametri le ceviane (segmenti che congiungono un vertice con il lato opposto). % lang-it
\end{proposition}

Bogdańska, Neugebauer \cite[s. 71]{neugebauer_2018} podają to jako ćwiczenie. % lang-pl
Bogdańska e Neugebauer \cite[p. 71]{neugebauer_2018} propongono questo risultato come esercizio. % lang-it

\begin{theorem}
[Auberta] % lang-pl
[di Aubert] % lang-it
	Dane są cztery proste w położeniu ogólnym (żadne trzy nie przechodzą przez jeden punkt, żadne dwie nie są równoległe). % lang-pl
	Wówczas ortocentra czterech trójkątów wyznaczonych przez trójki tych prostych leża na jednej prostej, zwanej prostą Auberta. % lang-pl
	\index{twierdzenie!Auberta} % lang-pl
	Siano date quattro rette in posizione generale (nessuna terna passa per uno stesso punto e nessuna coppia è parallela). % lang-it
	Allora gli ortocentri dei quattro triangoli determinati dalle terne di tali rette giacciono su una stessa retta, detta retta di Aubert. % lang-it
	\index{teorema!di Aubert} % lang-it
\end{theorem}

Nie udało nam się ustalić, kim będzie Aubert, ale najprawdopodobniej udowodni to twierdzenie w 1899 roku. % lang-pl
Non siamo riusciti a stabilire chi fosse Aubert, ma con ogni probabilità dimostrò questo teorema nel 1899. % lang-it
\index[persons]{Aubert, urodzony przed 1899 rokiem}%
Prosta Auberta znana jest czasami jako prosta Steinera. % lang-pl
La retta di Aubert è talvolta nota anche come retta di Steiner. % lang-it
% TODO: https://el-m-wikipedia-org.translate.goog/wiki/Ευθεία_Σίμσον_(τετράπλευρο)?_x_tr_sl=auto&_x_tr_tl=pl&_x_tr_hl=pl&_x_tr_pto=wapp
% https://el-m-wikipedia-org.translate.goog/wiki/Θεώρημα_Miquel_(τετράπλευρο)?_x_tr_sl=auto&_x_tr_tl=pl&_x_tr_hl=pl&_x_tr_pto=wapp
% https://el-m-wikipedia-org.translate.goog/wiki/Ευθεία_Νεύτωνα-Γκάους?_x_tr_sl=auto&_x_tr_tl=pl&_x_tr_hl=pl&_x_tr_pto=wapp
Eves \cite[s. 95]{eves1_1972} dopowiada, że trzy okręgi ze średnicami na przekątnych czworoboku zupełnego są współosiowe; środki trzech przekątnych leżą na prostej prostopadłej do Auberta. % lang-pl
Eves \cite[p. 95]{eves1_1972} aggiunge che le tre circonferenze aventi come diametri le diagonali di un quadrilatero completo sono coassiali; i punti medi delle tre diagonali giacciono su una retta perpendicolare alla retta di Aubert. % lang-it
% współosiowy = coaxial

\subsection{Pęki okręgów} % lang-pl
\subsection{Fasci di circonferenze} % lang-it

Pękiem (z angielskiego \emph{pencil}) będziemy nazywać dowolną rodziną geometrycznych obiektów (przede wszystkim prostych lub okręgów) o pewnej wspólnej cesze, którą można odtworzyć z każdej pary obiektów wspomnianej rodziny. % lang-pl
Chiameremo fascio (in inglese \emph{pencil}) una qualunque famiglia di oggetti geometrici (soprattutto rette o circonferenze) caratterizzata da una proprietà comune che può essere ricostruita a partire da una qualsiasi coppia di oggetti della famiglia. % lang-it

\begin{example}
	Do pęku prostych należą wszystkie proste, które przechodzą przez ustalony punkt. % lang-pl
	Appartengono a un fascio di rette tutte le rette che passano per un punto fissato. % lang-it
\end{example}

\begin{definition}
	Dowolne dwa niewspółśrodkowe okręgi $K_1, K_2$ wyznaczają jedną maksymalną rodzinę okręgów współosiowych z okręgami $K_1, K_2$, zwaną pękiem okręgów. % lang-pl
	Due qualunque circonferenze non concentriche $K_1$, $K_2$ determinano un'unica massima famiglia di circonferenze coassiali con $K_1$ e $K_2$, detta fascio di circonferenze. % lang-it
\end{definition}

O pękach okręgów pisze Neugebauer, Bogdańska \cite[s. 80]{neugebauer_2018}. % lang-pl
Dla ustalonch punktów $P, Q$ wprowadzają tam rodzinę % lang-pl
Dei fasci di circonferenze scrivono Neugebauer e Bogdańska \cite[p. 80]{neugebauer_2018}. % lang-it
Per punti fissati $P$, $Q$ introducono la famiglia % lang-it
\begin{equation}
	\mathcal H(PQ) := \{A(\lambda) : 0 \le \lambda \le + \infty, \lambda \neq 1 \},
\end{equation}
gdzie $A(\lambda)$ to okrąg Apoloniusza z definicji \ref{definicja_okregu_apoloniusza}. % lang-pl
Okręgi $A(\lambda)$ są parami rozłączne, wszystkie są prostopadłe do wszystkich okręgów przechodzących przez punkty $P$ i $Q$. % lang-pl
Uznają, że $A(\infty)$ to punkt $Q$, $A(0)$ to punkt $P$, zaś $A(1)$ to symetralna odcinka $PQ$. % lang-pl
Wtedy przez każdy punkt płaszczyzny przechodzi dokładnie jeden okrąg $A(\lambda)$. % lang-pl
dove $A(\lambda)$ è la circonferenza di Apollonio della definizione \ref{definicja_okregu_apoloniusza}. % lang-it
Le circonferenze $A(\lambda)$ sono a due a due disgiunte e tutte ortogonali a tutte le circonferenze passanti per i punti $P$ e $Q$. % lang-it
Assumono che $A(\infty)$ sia il punto $Q$, che $A(0)$ sia il punto $P$ e che $A(1)$ sia l'asse del segmento $PQ$. % lang-it
Allora per ogni punto del piano passa esattamente una circonferenza $A(\lambda)$. % lang-it

Maksymalną współosiową rodzinę okręgów parami rozłącznych nazywają pękiem hiperbolicznym okręgów, każdy jest postaci $\mathcal H(PQ)$ i punkty $P \neq Q$ są wyznaczone jednoznacznie. % lang-pl
Chiamano fascio iperbolico di circonferenze una massima famiglia coassiale di circonferenze a due a due disgiunte; ciascuna è della forma $\mathcal H(PQ)$ e i punti $P \neq Q$ sono determinati univocamente. % lang-it

\begin{definition}
	Punkty $P, Q$ nazywamy punktami Ponceleta, prostą $PQ$ linią środków, zaś symetralną odcinka $PQ$ -- osią pęku hiperbolicznego $\mathcal H(PQ)$. % lang-pl
	I punti $P$, $Q$ si chiamano punti di Poncelet, la retta $PQ$ si chiama linea dei centri e l'asse del segmento $PQ$ si chiama asse del fascio iperbolico $\mathcal H(PQ)$. % lang-it
\end{definition}

Nazwisko Ponceleta pojawia się tu nieprzypadkowo, ma związek z (małym) twierdzeniem Ponceleta. % lang-pl
Il nome di Poncelet non compare qui per caso: è legato al (piccolo) teorema di Poncelet. % lang-it

(Nie jest trudno się domyślić, że istnieją też pęki paraboliczne i eliptyczne). % lang-pl
(Non è difficile immaginare che esistano anche fasci parabolici ed ellittici). % lang-it

% en-Wiki: A pencil of circles (or coaxial system) is the set of all circles in the plane with the same radical axis.
% https://en.wikipedia.org/wiki/Pencil_(geometry) -> dowolne dwa obiekty pozwalają wyznaczyć wszystkie

Okręgi koncentryczne mają oś potęgową w nieskończoności. % lang-pl
Le circonferenze concentriche hanno l'asse radicale all'infinito. % lang-it
% pencil - pokok po rosyjsku, po polsku raczej pęk
O pęku piszą Audin \cite[s. 92-98]{audin_2003}, Eves \cite[s. 94]{eves1_1972}. % lang-pl
Dei fasci scrivono Audin \cite[p. 92--98]{audin_2003} ed Eves \cite[p. 94]{eves1_1972}. % lang-it

%